\documentclass[12pt,a4paper,french]{report}
% ─── PACKAGES STANDARDS ───────────────────────────────────────
\usepackage[utf8]{inputenc}
\usepackage[T1]{fontenc}
\usepackage{babel}
\usepackage{geometry}
\usepackage{xcolor}
\usepackage{hyperref}
\usepackage{graphicx}
\usepackage{listings}
\usepackage{fancyhdr}
\usepackage{titlesec}
\usepackage{tcolorbox}
\usepackage{enumitem}
\usepackage{booktabs}
\usepackage{longtable}
\usepackage{array}
\usepackage{tabularx}
\usepackage{float}
\usepackage{tikz}
\usepackage{amsmath}
\usepackage{setspace}
\usepackage{parskip}
\usepackage{charter}

% ─── GÉOMÉTRIE ───────────────────────────────────────────────
\geometry{
  top=2.5cm, bottom=2.5cm,
  left=2.5cm, right=2.5cm,
  headheight=15pt
}

% ─── COULEURS ────────────────────────────────────────────────
\definecolor{primaryblue}{RGB}{31,78,121}
\definecolor{lightblue}{RGB}{214,228,240}
\definecolor{accentblue}{RGB}{70,130,180}
\definecolor{darkgray}{RGB}{64,64,64}
\definecolor{lightgray}{RGB}{245,245,245}
\definecolor{successgreen}{RGB}{0,128,0}
\definecolor{warningorange}{RGB}{204,102,0}
\definecolor{errored}{RGB}{180,30,30}
\definecolor{codebg}{RGB}{248,248,248}

% ─── HYPERREF ────────────────────────────────────────────────
\hypersetup{
  colorlinks=true,
  linkcolor=primaryblue,
  urlcolor=accentblue,
  citecolor=primaryblue
}

% ─── EN-TÊTES / PIEDS DE PAGE ────────────────────────────────
\pagestyle{fancy}
\fancyhf{}
\fancyhead[L]{\small\color{darkgray} Panel Admin -- Documentation de Production}
\fancyhead[R]{\small\color{darkgray} \leftmark}
\fancyfoot[C]{\small\color{darkgray} \thepage}
\renewcommand{\headrulewidth}{0.4pt}
\renewcommand{\footrulewidth}{0pt}

% ─── TITRES ──────────────────────────────────────────────────
\makeatletter
\@ifundefined{chapter}{}{
  \titleformat{\chapter}[block]
    {\normalfont\huge\bfseries\color{primaryblue}}
    {\thechapter.}{1em}{}
    [\vspace{0.5ex}\titlerule[1pt]]
}
\makeatother

\titleformat{\section}[block]
  {\normalfont\Large\bfseries\color{primaryblue}}
  {\thesection}{1em}{}

\titleformat{\subsection}[block]
  {\normalfont\large\bfseries\color{accentblue}}
  {\thesubsection}{1em}{}

% ─── BOÎTES COLORÉES ─────────────────────────────────────────
\tcbuselibrary{skins,breakable}

\newtcolorbox{infobox}[1][]{
  enhanced, breakable,
  colback=lightblue!30, colframe=primaryblue,
  fonttitle=\bfseries\color{primaryblue},
  title={#1}, left=6pt, right=6pt, top=4pt, bottom=4pt
}

\newtcolorbox{warningbox}[1][Attention]{
  enhanced, breakable,
  colback=orange!10, colframe=warningorange,
  fonttitle=\bfseries\color{warningorange},
  title={#1}, left=6pt, right=6pt, top=4pt, bottom=4pt
}

\newtcolorbox{codebox}[1][]{
  enhanced, breakable,
  colback=codebg, colframe=darkgray!40,
  fonttitle=\bfseries\color{darkgray},
  title={#1}, left=6pt, right=6pt, top=4pt, bottom=4pt
}

% ─── COMMANDES PERSONNALISÉES ────────────────────────────────
\newcommand{\bucket}[1]{\texttt{s3://#1}}
\newcommand{\scriptname}[1]{\texttt{#1}}
\newcommand{\ymlkey}[1]{\textbf{\texttt{#1}}}


\title{Architecture Technique : Pipeline Panel Admin}
\author{Data Engineering Team}
\date{\today}

\begin{document}

\maketitle
\tableofcontents
\newpage

\chapter{Architecture Globale et Voyage de la Donnée}

Le système repose sur une architecture \textbf{Medallion (ou Data Lakehouse)}, permettant une montée en maturité progressive de la donnée. Ce choix garantit que les transformations métier sont isolées des opérations de nettoyage technique.

\section{Le "Data Journey" : de l'Excel au Parquet Gold}

\begin{enumerate}
    \item \textbf{Ingestion (Staging)} : Les fichiers Excel mensuels sont déposés sur MinIO. Cette zone est la seule contenant des formats non-optimisés et hétérogènes (noms de colonnes variables, en-têtes décalés).
    \item \textbf{Uniformisation (Bronze)} : Chaque Excel est converti en Parquet. Toutes les colonnes sont conservées en type \textit{String} pour éviter toute perte d'information. C'est ici qu'on réalise la jointure technique F1-F2 sur la clé \texttt{matricule||organisme}.
    \item \textbf{Métier et Normalisation (Silver)} : Les données sont typées. Les nomenclatures administratives (ANSTAT) sont rattachées via un \textbf{fuzzy join} textuel. Les agents hors-périmètre (situations non-actives) sont supprimés. Les fichiers sont regroupés par année civile.
    \item \textbf{Stabilisation et Imputation (Gold)} : Le panel longitudinal est constitué. Les grades manquants sont imputés par cascade statistique. Cette couche est la "Source de Vérité" pour toutes les analyses.
\end{enumerate}

\section{Stratégie de Stockage : MinIO et Parquet}

Le projet utilise \textbf{MinIO}, un stockage objet S3-compatible. Cela permet de simuler un environnement Cloud moderne tout en restant sur une infrastructure locale souveraine.
Le format \textbf{Apache Parquet} offre des lectures colonnaire extrêmement rapides et une compression binaire réduisant l'espace disque de 90\% par rapport à l'Excel.

\chapter{Dictionnaire des Données Gold}

Le panel final (\scriptname{panel\_complet.parquet}) contient les colonnes suivantes :

\begin{longtable}{lll}
  \toprule
  \textbf{Groupe} & \textbf{Colonne} & \textbf{Type} \\
  \midrule
  \endhead
  Identité & \texttt{matricule} & Utf8 \\
  & \texttt{nom} & Utf8 \\
  & \texttt{sexe} & Utf8 \\
  & \texttt{date\_naissance} & Utf8 \\
  \midrule
  Nomenclature & \texttt{CODE\_ORGANISME} & Utf8 (Code ANSTAT) \\
  & \texttt{organisme} & Utf8 (Libellé source) \\
  & \texttt{CODE\_EMPLOI} & Utf8 \\
  & \texttt{Metier\_CITP} & Utf8 \\
  & \texttt{GRADE} & Utf8 (Imputé/Observé) \\
  & \texttt{GRADE\_SOURCE} & Utf8 (OBSERVE|IMPUTE|NF) \\
  \midrule
  Temporel & \texttt{mois\_annee} & String (Format MMYYYY) \\
  & \texttt{annee} / \texttt{mois} & Int32 \\
  \midrule
  Salaires & \texttt{montant\_brut} & Float64 (F1) \\
  & \texttt{montant\_net} & Float64 (F1) \\
  & \texttt{124\_f2} \ldots \texttt{425\_f2} & Float64 (Composantes détaillées) \\
  \midrule
  Traçabilité & \texttt{fichier\_source} & Utf8 \\
  \bottomrule
\end{longtable}

\chapter{Moteur d'Exécution et Performance}

\section{Streaming et Gestion de la RAM avec Polars}

Le pipeline utilise \textbf{Polars}. Pour traiter des millions de lignes sans saturer la RAM, nous exploitons le \textbf{Lazy Evaluation} :
\begin{itemize}
    \item \texttt{scan\_parquet()} : Lit uniquement les métadonnées pour préparer le plan d'exécution.
    \item \texttt{sink\_parquet()} : Écrit le résultat en streaming vers le disque, permettant de traiter des volumes supérieurs à la RAM disponible.
\end{itemize}

\section{Configuration Dynamique (YAML)}

La logique métier est séparée du code Python via \ymlkey{pipeline\_settings.yaml} :
\begin{itemize}
    \item \ymlkey{column\_mapping} : Gère la variabilité des en-têtes Excel via regex.
    \item \ymlkey{situations\_valides} : Définit le périmètre d'activité administrative.
\end{itemize}

\section{Observabilité Industrielle}

Le module \scriptname{common.logging\_config} remplace les \texttt{print()}. Il assure une traçabilité horodatée et structurée, permettant à l'orchestrateur de générer des rapports de run Markdown détaillés pour chaque exécution.

\end{document}
